\subsection{Integral leading minors and their arithmetic sum}
\label{shared-integral-leading-minors}

The arithmetic normalization of the leading minors is useful before
the sum over rational returns is taken.  The prescribed weighted flag
lattices give more than integrality of the unnormalized determinants:
they give their exact different-ideal divisibility.  We prove that
divisibility directly, use it to sum the actual labels, and then
combine the real-place orbital estimates.  The estimate for the
number of returns with a fixed label remains a separate arithmetic
requirement.  All source identifications below refer to the pinned
revision in~\cite{global-gram-integral-minors-source}.

\paragraph{The original weighted lattices.}
Let $E/F$ be the original CM extension, let $[F:\mathbf Q]=f$,
and write $\mathcal O=\mathcal O_E$.  Complex conjugation is denoted
by a bar.  Let $\mathfrak d=\mathfrak D_{F/\mathbf Q}\mathcal O$
be the extension of the different of $F$.  It is a nonzero integral
ideal fixed by conjugation.  In the original family $E/F$ is
unramified at finite places.  Consequently, if
$\rho=|\operatorname{disc} E|^{1/(2f)}$, then
\begin{equation}\label{eq:minor-different-norm}
 N\mathfrak d=|\operatorname{disc}F|^2
             =|\operatorname{disc}E|=\rho^{2f}.
\end{equation}
The first equality follows by extending an ideal through a degree-two
extension, and the second is the relative discriminant formula.

Let $V$ carry the original Hermitian form $h$, conjugate-linear in
its first variable, and let $\Lambda\subset V$ be the integral
lattice common to the two flags.  Its original unimodularity
implies $h(\Lambda,\Lambda)\subset\mathcal O$.
Number the isotropic vectors in the flags as
$e_0,\ldots,e_{r-1}$ and $e'_0,\ldots,e'_{r-1}$.
The prescribed flag lattice is
\[
 X_{\mathcal O}(e)=\sum_{a=0}^{r-1}\mathfrak d^{-(r-a)}e_a.
\]
The original two flag conditions say
$X_{\mathcal O}(e)\subset\Lambda$ and
$X_{\mathcal O}(e')\subset\Lambda$; in fact each is the indicated
intersection of $\Lambda$ with its isotropic subspace.
For $0\le l\le r$, put
\[
 a_l=\sum_{a=0}^{l-1}(r-a),\qquad
 \delta_l=\det\bigl(h(e'_b,e_a)\bigr)_{0\le a,b<l},
 \qquad \delta_0=1.
\]
Changing the convention for which flag is listed first conjugates a
minor and leaves every ideal norm and estimate below unchanged.

\begin{theorem}[Divisibility of the prescribed minors]
\label{thm:actual-minor-ideal}
For every $0\le l\le r$,
\[
 \delta_l\in\mathfrak d^{2a_l}.
\]
If $\delta_l\ne0$, the fractional ideal
$\mathfrak a_l=(\delta_l)\mathfrak d^{-2a_l}$ is a nonzero
integral ideal.  In particular,
\begin{equation}\label{eq:actual-minor-product}
 1\le N\mathfrak a_l
   =|N_{E/\mathbf Q}\delta_l|(N\mathfrak d)^{-2a_l}
   =\prod_{\tau\in\Phi}
       \bigl(|\tau(\delta_l)|\rho^{-2a_l}\bigr)^2,
\end{equation}
where $\Phi$ contains one embedding from each conjugate pair.
\end{theorem}

\begin{proof}
For $x\in\mathfrak d^{-(r-a)}$ and
$y\in\mathfrak d^{-(r-b)}$, both $xe_a$ and $ye'_b$ belong
to $\Lambda$.  Integrality of $h$ gives
$\overline y x h(e'_b,e_a)\in\mathcal O$.
Since $\overline{\mathfrak d}=\mathfrak d$ and every nonzero
fractional ideal of $\mathcal O$ is invertible, this is exactly
\[
 h(e'_b,e_a)\in\mathfrak d^{2r-a-b}.
\]
For example, the assertion follows after localization at each prime
ideal, where multiplying by the two inverse ideal powers subtracts
their valuations.  Each summand in the determinant expansion has
ideal valuation at least
$\sum_{a=0}^{l-1}(2r-a-\sigma(a))=2a_l$.
This proves the divisibility without choosing a generator of
$\mathfrak d$ or a basis of $\Lambda$ over $\mathcal O$.
The ideal $\mathfrak a_l$ is therefore integral.  Its norm is a
positive integer if $\delta_l\ne0$.  The last identity follows
from~\eqref{eq:minor-different-norm} and the product formula for
the field norm in a CM field.
\end{proof}

This proves the normalized lower bound directly for every nonzero
leading minor of the original flags.  Unnormalized integrality
alone would not do so: dividing the norm of an arbitrary integral
element by $(N\mathfrak d)^{2a_l}$ can give a number smaller than
one.  The weighted flag condition is essential.

\paragraph{A packing estimate with the actual normalization.}
We use the following elementary estimate.  For a finite set
$P\subset\mathbf C^f$, suppose
\[
 \prod_{\tau=1}^f|p_\tau-q_\tau|\ge1
       \quad(p,q\in P,\ p\ne q).
\]
For $R\ge1$ one then has
\begin{equation}\label{eq:minor-packing}
 \sum_{p\in P}\prod_{\tau=1}^f
        \max\{1,|p_\tau|/R\}^{-3}\le(64R)^{2f}.
\end{equation}
Here is a proof retaining an absolute constant.  Assign $z\in\mathbf C$
to the square with indices
$(\lfloor2\operatorname{Re}z\rfloor,
  \lfloor2\operatorname{Im}z\rfloor)$.
Two tuples in the same product of squares have every coordinate
distance smaller than $1$, so at most one point of $P$ occupies
each such product.  Set
\[
 w(n)=\max\{1,(|n|-1)/(2R)\}^{-3/2}\quad(n\in\mathbf Z).
\]
The weight of $z$ is at most the product of the two corresponding
$w$'s.  Indeed $|z|$ is at least each real coordinate's absolute
value, and
$\max(1,|z|/R)^2$ is at least the product of the two coordinate
maxima.  The floor inequalities then give the claimed bound.
The decreasing function $u\mapsto\max(1,u/(2R))^{-3/2}$ has
integral $6R$ on $[0,\infty)$.  Hence
\[
 \sum_{n\in\mathbf Z}w(n)
 =3+2\sum_{n=1}^{\infty}\max\{1,n/(2R)\}^{-3/2}
 \le3+12R\le15R.
\]
Summing over all product squares proves~\eqref{eq:minor-packing},
in fact with $15$ in place of $64$.

Fix $H\ge1$, $1\le r$, and $0\le j\le r$.  For a label
$\delta=(\delta_0,\delta_1,\ldots)$ with
$\delta_0=1$, $\delta_l=0$ for $l>j$, and
$\delta_l\in\mathfrak d^{2a_l}$ for $1\le l\le j$, define
\begin{align*}
 D_{\tau l}(\delta)&=|\tau(\delta_l)|\rho^{-2a_l},&
 M&=(2H)^{2r},\\
 \sigma_\tau(\delta)&=
    \max\bigl(1,D_{\tau l}(\delta)/M:1\le l\le j\bigr),&
 S(\delta)&=\prod_{\tau\in\Phi}\sigma_\tau(\delta).
\end{align*}
Zero internal minors are permitted.  They contribute zero to the
maximum; they are never inverted.

\begin{theorem}[The actual leading-minor sum]
\label{thm:actual-normalized-minor-sum}
For any set $\mathcal T$ of distinct labels satisfying these
conditions, the nonnegative series is summable and
\begin{equation}\label{eq:normalized-label-sum}
 \sum_{\delta\in\mathcal T}S(\delta)^{-3r}
       \le2^j\bigl(64(2H)^{2r}\bigr)^{2fj}.
\end{equation}
There is no factor depending on the root discriminant in this
bound.  For the original scale $H=b^{C_Hm}$, with $b\ge2$,
$m=2r+2$ and $C_H\ge0$, it implies the original consumer's bound
\begin{equation}\label{eq:original-minor-label-budget}
 \sum_{\delta\in\mathcal T}S(\delta)^{-3r}
     \le b^{C_c^2m^2fr},\qquad C_c=21+4C_H.
\end{equation}
The constants are chosen before the field, rank, words, and finite
table entries.
\end{theorem}

\begin{proof}
For a fixed $l$, consider the normalized embedding
$x\mapsto(\rho^{-2a_l}\tau(x))_{\tau\in\Phi}$ of
$\mathfrak d^{2a_l}$.  If $x\ne y$, their difference is a
nonzero element of the same ideal.  The proof of
Theorem~\ref{thm:actual-minor-ideal} applied to this element gives
\[
 \prod_{\tau\in\Phi}
       |\rho^{-2a_l}\tau(x-y)|\ge1.
\]
Thus~\eqref{eq:minor-packing} applies to every finite subset of
this exact normalized ideal, with $R=M$.

Partition any finite subset of $\mathcal T$ according to
$A=\{1\le l\le j:\delta_l\ne0\}$.
Since $|A|\le r$ and $\sigma_\tau\ge1$,
\[
 S(\delta)^{-3r}\le
 \prod_{l\in A}\prod_{\tau\in\Phi}
       \max\{1,D_{\tau l}(\delta)/M\}^{-3}.
\]
For fixed $A$ the selected coordinates determine the label.
Enlarge their allowed sets to the product of their individual
coordinate sets and apply~\eqref{eq:minor-packing} separately to
each ideal.  The resulting sum is at most $(64M)^{2f|A|}$.
There are $2^j$ possible supports, so the finite sum is bounded
by the right side of~\eqref{eq:normalized-label-sum}.
The labels form a countable set; increasing finite subsets proves
both summability and the same bound for the full series.
For $j=0$ there is one possible label and the bound is $1$.

For the final assertion take logarithms to base $b$.  The right
side of~\eqref{eq:normalized-label-sum} has logarithm at most
\[
 j+2fj\{6+2r(1+C_Hm)\}
 \le fr\{13+4r+4C_Hmr\}
 \le (17+4C_H)m^2fr
 \le C_c^2m^2fr.
\]
Here $f,r\ge1$, $j\le r\le m$ and $b\ge2$ were the only
inequalities used.  In particular this step does not need an
extra restriction relating the root discriminant to $b$.
\end{proof}

\paragraph{The normalized discriminant bound.}
For every nonzero selected minor define
$\Delta_l=N((\delta_l)\mathfrak d^{-2a_l})$.
Equations~\eqref{eq:actual-minor-product} and the definition of
$\sigma_\tau$ give
\begin{equation}\label{eq:minor-height-bound}
 1\le\Delta_l\le M^{2f}S(\delta)^2
       \le b^{C m^2f}S(\delta)^2
       \quad\bigl(C\ge2(C_H+1)\bigr).
\end{equation}
Indeed $\log_bM^{2f}\le4rf(1+C_Hm)\le2(C_H+1)m^2f$.
The product of the nonzero minor norms is correspondingly bounded by
$(b^{Cm^2f}S^2)^r$.  These are the original normalized ideal norms,
not the norms of unnormalized integral determinants.

\paragraph{Combining the real-place estimates.}
For one actual relevant nonidentity return with zero Gram
determinant, let $1\le j\le r$ be its original index.  The
original normal form has $\delta_j\ne0$ and
$\delta_l=0$ for $j<l\le r$.  Set $\Delta=\Delta_j$ and
$a=k-1-j$.  The local minor dictionary identifies $D_{\tau l}$
above with the leading-minor absolute values in
Theorem~\ref{thm:actual-local-degenerate-bound}.
Retain the original positive Haar factors $\kappa_\tau/s_\tau$.

\begin{theorem}[The product of the original local bounds]
\label{thm:actual-degenerate-product}
Suppose $k\ge40r^4$ and $k\ge5r$.  Fix $C_s\ge0$ and
$q=(\lceil C_s\rceil+2)kr$.  Under the original local hypotheses
and their proved moment allowance, the actual local orbital
operators satisfy
\begin{align}
 \prod_{\tau\in\Phi}\|J_\tau\|_{\mathrm{op}}
 &\le\left(\prod_{\tau\in\Phi}\frac{\kappa_\tau}{s_\tau}\right)
       \left(\frac{16\pi e}{k}\right)^{j(k-1-j)f}
       \Delta^{-(k-1-j)/2}\notag\\
 &\hspace{12mm}\cdot(8rH)^{80r^4f}S(\delta)^{-C_sk}.
 \label{eq:actual-degenerate-product}
\end{align}
Every factor on the right has the original arithmetic
normalization.  The finite-place stabilizer factors have not
been cancelled or discarded.
\end{theorem}

\begin{proof}
Put $E_0=8r^2+2$ and
$B_\tau=\max(1,D_{\tau l}^{-1}:0\le l\le r,\ D_{\tau l}>0)$.
Theorem~\ref{thm:actual-minor-ideal} says that a nonzero minor
has positive values at all places and their product is at least
one.  For each such $l$,
\[
 \prod_\tau\max(1,D_{\tau l}^{-1})
 =\frac{\prod_\tau\max(1,D_{\tau l})}{\prod_\tau D_{\tau l}}
 \le\prod_\tau\max(1,D_{\tau l}).
\]
Using $\max(1,D_{\tau l})\le M\sigma_\tau$, including
$l=0$, gives
\begin{equation}\label{eq:reciprocal-minor-product}
 \prod_\tau B_\tau\le M^{f(r+1)}S^{r+1}.
\end{equation}
All real-place Haar scalars remain unchanged while we multiply
Theorem~\ref{thm:actual-local-degenerate-bound}.
The powers of $D_{\tau j}$ combine to $\Delta^{-a/2}$.
Since $a\ge k/2$, its factor $8\pi e/a$ is at most
$16\pi e/k$.

The total exponent of $S$ after
\eqref{eq:reciprocal-minor-product} is
\[
 (r+1)(8r^2+2)-\frac{q-16r^3-4r-1}{r}
 \le-C_sk.
\]
To check this, the positive terms are at most
$8r^3+24r^2+2r+6+1/r\le41r^3\le2k$,
whereas $q/r\ge(C_s+2)k$.
The remaining constant at each place is
\[
 16(4r+4)(4r)^{4r^2}
 (2H)^{16r^3+4r+2r(r+1)(8r^2+2)}.
\]
Since $16(4r+4)\le(8r)^3$ and $H\ge1$, it is at most
$(8rH)^{16r^4+32r^3+8r^2+8r+3}$, hence at most
$(8rH)^{80r^4}$.  This proves the stated product estimate.
The chosen $q$ exceeds $16r^3+4r+2$ under $k\ge40r^4$;
the original moment theorem allows the fixed constant
$\lceil C_s\rceil+2$ before the arithmetic data are chosen.
\end{proof}

\paragraph{What remains in the arithmetic sum.}
The preceding theorems supply the exact minor-label sum, the
normalized discriminant upper bound and the product of local
orbital bounds.  They apply to the actual weighted flag lattices
in both the RH and Dirichlet-GRH routes.  The next requirement is
a count of the returns sharing each label, with the original
inverse-automorphism weights and stabilizer covolumes.
Those multiplicities are absent from~\eqref{eq:normalized-label-sum}.
The source obtains them by a mean over a genus of definite
Hermitian lattices, including pairs of embedded data and their
isometries.  A count of distinct minors does not establish that
mean.  Until it is proved and combined with the remaining
normalizers, neither the full zero-Gram operator bound nor the
final positive trace estimates follow from this subsection.
